%% LyX 2.4.4 created this file.  For more info, see https://www.lyx.org/.
%% Do not edit unless you really know what you are doing.
\documentclass[oneside,english]{amsart}
\usepackage[T1]{fontenc}
\usepackage[latin9]{inputenc}
\synctex=-1
\usepackage{mathrsfs}
\usepackage{enumitem}
\usepackage{amsthm}

\makeatletter
%%%%%%%%%%%%%%%%%%%%%%%%%%%%%% Textclass specific LaTeX commands.
\newlength{\lyxlabelwidth}      % auxiliary length 
\theoremstyle{definition}
\newtheorem*{defn*}{\protect\definitionname}
\theoremstyle{plain}
\newtheorem{thm}{\protect\theoremname}
\theoremstyle{plain}
\newtheorem{prop}[thm]{\protect\propositionname}

%%%%%%%%%%%%%%%%%%%%%%%%%%%%%% User specified LaTeX commands.
\date{}
\usepackage{commath}

\makeatother

\usepackage{babel}
\providecommand{\definitionname}{Definition}
\providecommand{\propositionname}{Proposition}
\providecommand{\theoremname}{Theorem}

\begin{document}
\begin{defn*}
A lattice is \emph{proper} if none of its cells is either $\begin{array}{cc}
0 & 1\\
1 & 0
\end{array}$ or $\begin{array}{cc}
1 & 0\\
0 & 1
\end{array}$.
\end{defn*}
%
\begin{defn*}
Let $P$ be a polygon mosaic. If we replace blank tiles of $P$ (all,
some, or none) with other tiles, then the resulting mosaic $M$ is
called a \emph{mosaic created from $P$}.
\end{defn*}
\begin{prop}
For a polygon mosaic $P$, let $\ell=\mathscr{L}\left(P\right)$.
\begin{enumerate}
\item The number of mosaics created from $P$ equals $\left|v\left(\ell\right)\right|$.
\item If $P$ contains an even {[}resp. odd{]} number of polygons then $v\left(\ell\right)$
is positive {[}resp. negative{]}.
\end{enumerate}
\end{prop}

\begin{proof}
The first part is simple. Observe that $\left|v\left(\ell\right)\right|=7^{i}$,
where $i$ is the number of blank tiles of $P$, and this is equal
to the number of mosaics created from $P$. It is in the second part
where we will need a lot more work.

Let $N$ be the number of vertices of $\ell$ that are $1$, and define
$\ell_{0}$ to be the lattice formed by changing these $N$ vertices
from $1$ to $0$. That is, $\ell_{0}$ has the same dimensions as
$\ell$ but all vertices are $0$. We will change $\ell_{0}$ back
to $\ell$ by changing these vertices back from $0$ to $1$, one
vertex at a time. The order in which we make these changes is as follows.

We change back the vertices one row at a time, working from top to
bottom. In each row we make two passes, each time moving from left
to right, and each time only changing vertices if the corresponding
vertex of $\ell$ is $1$. On the first pass we change a vertex if
the vertex is directly above it is $1$. (Because we are working the
rows from top to bottom, this vertex is $1$ in both our current lattice
and also in $\ell$.) On the second pass we change the remaining vertices
of the row, i.e., if the vertex is directly above it is $0$.

This procedure creates a sequence $\ell_{0}$, $\ell_{1}$, \ldots ,
$\ell_{N}$ of lattices with $\ell_{N}=\ell$. Because $\ell=\mathscr{L}\left(P\right)$,
we know $\ell$ is a proper framed lattice. Since every $\ell_{i}$
has the same dimensions as $\ell$ and has vertices of $0$ in at
least the same positions as in $\ell$, we see every $\ell_{i}$ is
a framed lattice. We will prove that every $\ell_{i}$ is also a proper
lattice (which is the reason we make two passes through each row)
and that by taking $P_{i}=\mathscr{L}^{-1}\left(\ell_{i}\right)$,
the second part of the proposition holds for $P_{i}$ and $\ell_{i}$.

Because every $\ell_{i}$ is framed, in the transition from $\ell_{i-1}$
to $\ell_{i}$, the single vertex that changes is adjacent to eight
vertices, as in Figure~TBA (please make it look like $\begin{array}{ccc}
ul & u & ur\\
l & \# & r\\
dl & d & dr
\end{array}$), where the changing vertex $\#$ is $0$ in $\ell_{i-1}$ and $1$
in $\ell_{i}$, and the remaining vertices are labeled for their relative
position to this central vertex, such as $dl$ for down left and $u$
for up.

We need to consider all possible adjacent vertices. Since our process
goes from top to bottom one row at a time, we have $\left(dl,d,dr\right)=\left(0,0,0\right)$.
Also, because of the way in which we make our two passes, if $u=1$
then $r=0$ and $\left(ul,l\right)\neq\left(0,1\right)$, whereas
if $u=0$ then $\left(ur,r\right)\neq\left(0,1\right)$. Additionally,
suppose $u=0$, and that neither $\left(ul,l\right)$ nor $\left(ur,r\right)$
equals $\left(1,0\right)$. Indeed, if $\left(ul,l\right)=\left(1,0\right)$
then since we are making the second pass through the row, we would
have a cell $\begin{array}{cc}
1 & 0\\
0 & 1
\end{array}$ in $\ell_{i}$ and thus also in $\ell_{N}=\ell$, contradicting that
$\ell$ is a proper lattice. Similarly, $\left(ur,r\right)=\left(1,0\right)$
implies the that $\ell$ has a cell $\begin{array}{cc}
0 & 1\\
1 & 0
\end{array}$ (since with $ur=1$, $r$ cannot change in the second pass), again
contradicting that $\ell$ is a proper lattice. Putting all this together,
on the first pass, i.e., when $u=1$, the vertices can only be as
in the six cases pictured in Figure~{[}5{]}. Similarly, on the second
pass we have $u=0$ and one of the six cases of Figure~{[}7{]}.

Since trivially $\ell_{0}$ is a proper lattice, the twelve cases
in Figures~{[}5{]} and {[}7{]} show inductively that every $\ell_{i}$
is a proper lattice. For every $i$, we have already seen that $\ell_{i}$
is a framed lattice, so let $P_{i}=\mathscr{L}^{-1}\left(\ell_{i}\right)$.
Figures~{[}5{]} and {[}7{]} also show the portions of $P_{i-1}$
and $P_{i}$ that are in these two cells, where the dashed(?) lines
are from the polygons of $P_{i-1}$ and the dotted(?) lines from the
polygons of $P_{i}$. (Recall that $\#$ is $0$ in $\ell_{i-1}$
and $1$ in $\ell_{i}$.)

Let us say a cell is \emph{negative} if it is $\begin{array}{cc}
0 & 0\\
1 & 0
\end{array}$ or $\begin{array}{cc}
1 & 0\\
1 & 1
\end{array}$, i.e., if its evaluation under $v$ is negative. The second conclusion
of the proposition can be rephrased by saying that the number of negative
cells in $\ell$ and the number of polygons in $P$ have the same
parity, i.e., are both even or both odd. For example, $\ell_{0}$
satisfies the conclusion of the proposition because $\ell_{0}$ has
no negative cells and $P_{0}$ has no polygons.

Assume now that some $P_{i-1}$ and $\ell_{i-1}$ satisfies the second
conclusion of the proposition, and let us show that the same is true
for $P_{i}$ and $\ell_{i}$. This will complete the proposition.
We know that we are in one of the twelve cases of Figures~{[}5{]}
and {[}7{]}, and we must look at the number of negative cells in $\ell_{i}$
versus $\ell_{i-1}$, and the number of polygons in $P_{i}$ versus
$P_{i-1}$. An easy case is (a) of Figure~{[}7{]} because $\ell_{i}$
has one more negative cell (specifically, a cell $\begin{array}{cc}
0 & 0\\
1 & 0
\end{array}$), and $P_{i}$ has one more polygon than $P_{i-1}$.

Let us now look in the remaining cases at how the number of negative
cells changes from $\ell_{i-1}$ to $\ell_{i}$. In the six cases
of Figure~{[}5{]} and also case (d) of Figure~{[}7{]}, the four
pictured cells have no negative cells regardless of whether $\#$
equals $0$ or $1$, so there is no change. The number is also unchanged
in (b) of Figure~{[}7{]} because in this case a single negative cell
$\begin{array}{cc}
0 & 0\\
1 & 0
\end{array}$ changes position between $\ell_{i-1}$ and $\ell_{i}$. In (c) of
Figure~{[}7{]}, the number of negative cells increases by $2$ because
changing $0$ to $1$ in the central position creates both a cell
$\begin{array}{cc}
1 & 0\\
1 & 1
\end{array}$ and a cell $\begin{array}{cc}
0 & 0\\
1 & 0
\end{array}$. For the final two cases, the number decreases by $1$ in (e) of
Figure~{[}7{]} and increases by $1$ in (f) of Figure~{[}7{]}. Therefore,
we can complete the proposition by showing that in the six cases of
Figure~{[}5{]} and in (b), (c), (d) of Figure~{[}7{]}, $P_{i-1}$
and $P_{i}$ have the same number of polygons, whereas in (e) and
(f) of Figure~{[}7{]}, the number of polygons in $P_{i-1}$ and $P_{i}$
differ by $1$.

Consider, for example, (c) of Figure~{[}5{]}, and look at the point
just above the lower left corner and the point just left of the upper
right corner. In $P_{i-1}$ these are connected by a straight dashed
line, and in $P_{i}$ these are connected by dotted line segments.
We call both of these \emph{inside connections}. Because $P_{i-1}$
and $P_{i}$ are polygon mosaics, these same points are connected
by straight line segments outside the pictured area, and we will call
these \emph{outside connections} because they connect points on the
edge of the pictured square with a path outside the square. This outside
connection is the same in $P_{i-1}$ and $P_{i}$. Therefore $P_{i-1}$
and $P_{i}$ have the same number of polygons because the polygon
with this outside connection changes shape, and all other polygons
are unchanged. Similar reasoning completes the remaining cases of
Figure~{[}5{]} and cases (b), (c), (d) of Figure~{[}7{]}.

We now must show that in (e) and (f) of Figure~{[}7{]}, the number
of polygons in $P_{i-1}$ and $P_{i}$ differ by $1$. We will explicitly
show this in (e), and essentially identical reasoning applies to (f).
There are four places where inside connections meet the edge of the
outer square. Beginning in the lower left and working clockwise, let
$x_{1}$, $x_{2}$, $x_{3}$, and $x_{4}$ be these points. For example,
$x_{3}$ is just left of the upper right corner and $x_{4}$ is just
above the lower right corner. In $P_{i-1}$, there is an inside connection
between $x_{1}$ and $x_{2}$, and another between $x_{3}$ and $x_{4}$,
whereas $P_{i}$ has inside connections between $x_{1}$ and $x_{4}$,
and between $x_{2}$ and $x_{3}$.

Any outside connections in $P_{i-1}$ and $P_{i}$ must be identical,
and since these are polygon mosaics, each of $x_{1}$, $x_{2}$, $x_{3}$,
and $x_{4}$ is the end of some outside connection. One possibility
is that we have one outside connection between $x_{1}$ and $x_{2}$,
and another between $x_{3}$ and $x_{4}$. In this case, the inside
and outside connections of $P_{i-1}$ complete two polygons, one including
$x_{1}$ and $x_{2}$, and the other including $x_{3}$ and $x_{4}$.
However, $P_{i}$ will have one fewer polygon because a single polygon
contains all $x_{i}$. Another possibility is to have one outside
connection between $x_{1}$ and $x_{4}$, and another between $x_{2}$
and $x_{3}$. In this case, in $P_{i-1}$ a single polygon includes
all $x_{i}$, whereas in $P_{i}$ they complete two polygons, one
containing $x_{1}$ and $x_{4}$, and the other containing $x_{2}$
and $x_{3}$, and therefore $P_{i}$ has one more polygon than $P_{i-1}$.
These are the only possibilities since it is impossible to have an
outside connection between $x_{1}$ and $x_{3}$, and another between
$x_{2}$ and $x_{4}$. This is because our polygons are in a plane,
and polygons in a mosaic are not able to cross one another. This completes
our proof.
\end{proof}

\end{document}
